
\documentclass{article}
\usepackage{colortbl}
\usepackage{makecell}
\usepackage{multirow}
\usepackage{supertabular}

\begin{document}

\newcounter{utterance}

\centering \large Interaction Transcript for game `hot\_air\_balloon', experiment `air\_balloon\_survival\_en\_reasoning off\_hard', episode 5 with deepseek{-}chat{-}v3.1{-}t1.0.
\vspace{24pt}

{ \footnotesize  \setcounter{utterance}{1}
\setlength{\tabcolsep}{0pt}
\begin{supertabular}{c@{$\;$}|p{.15\linewidth}@{}p{.15\linewidth}p{.15\linewidth}p{.15\linewidth}p{.15\linewidth}p{.15\linewidth}}
    \# & $\;$A & \multicolumn{4}{c}{Game Master} & $\;\:$B\\
    \hline

    \theutterance \stepcounter{utterance}  
    & & \multicolumn{4}{p{0.6\linewidth}}{
        \cellcolor[rgb]{0.9,0.9,0.9}{
            \makecell[{{p{\linewidth}}}]{
                \texttt{\tiny{[P1$\langle$GM]}}
                \texttt{You are participating in a collaborative negotiation game.} \\
\\ 
\texttt{Together with another participant, you must agree on a single set of items that will be kept. Each of you has your own view of how much each item matters to you (importance). You do not know how the other participant values the items. Additionally, you are given the effort each item demands.} \\
\texttt{You may only agree on a set if the total effort of the selected items does not exceed a shared limit:} \\
\\ 
\texttt{LIMIT = 2307} \\
\\ 
\texttt{Here are the individual item effort values:} \\
\\ 
\texttt{Item effort = \{"A21": 147, "B93": 698, "A11": 588, "B33": 43, "C08": 356, "B39": 80, "B96": 95, "C85": 743, "A82": 106, "B31": 308, "B20": 325, "A60": 255, "B09": 276, "C81": 543, "B21": 51\}} \\
\\ 
\texttt{Here is your personal view on the importance of each item:} \\
\\ 
\texttt{Item importance values = \{"A21": 138, "B93": 583, "A11": 65, "B33": 262, "C08": 121, "B39": 508, "B96": 461, "C85": 484, "A82": 668, "B31": 389, "B20": 215, "A60": 97, "B09": 500, "C81": 30, "B21": 400\}} \\
\\ 
\texttt{Goal:} \\
\\ 
\texttt{Your goal is to negotiate a shared set of items that benefits you as much as possible (i.e., maximizes total importance to YOU), while staying within the LIMIT. You are not required to make a PROPOSAL in every message {-} you can simply negotiate as well. All tactics are allowed!} \\
\\ 
\texttt{Interaction Protocol:} \\
\\ 
\texttt{You may only use the following structured formats in a message:} \\
\\ 
\texttt{PROPOSAL: \{'A', 'B', 'C', …\}} \\
\texttt{Propose keeping exactly those items.} \\
\texttt{REFUSE: \{'A', 'B', 'C', …\}} \\
\texttt{Explicitly reject opponent's proposal.} \\
\texttt{ARGUMENT: \{'...'\}} \\
\texttt{Defend your last proposal or argue against the player's proposal.} \\
\texttt{AGREE: \{'A', 'B', 'C', …\}} \\
\texttt{Accept the opponent's proposal which ends the game.} \\
\\ 
\\ 
\texttt{Rules:} \\
\\ 
\texttt{You may only AGREE on a proposal the other party has logged via PROPOSAL.} \\
\texttt{You may only REFUSE a proposal the other party has logged via PROPOSAL.} \\
\texttt{Total effort of any PROPOSAL or AGREE set must be ≤ LIMIT.} \\
\texttt{Do NOT reveal your hidden importance scores.} \\
\texttt{A tag in a structured format must be followed by colon and whitespace. The argument must be a python set containing 0 or more strings.} \\
\texttt{So, it must be of the form TAG: \{...\}} \\
\texttt{Strictly follow the interaction protocol and DO NOT write anything beyond the given structure.} \\
\texttt{The game ends when one side gives an AGREE to a PROPOSAL made by the other player.} \\
\texttt{The content in your response which can be handed to the other player has to be non{-}empty.} \\
\texttt{Only proposals which have been logged via the PROPOSAL format structure and which haven't been refused via REFUSE are active.} \\
\texttt{You must include the ARGUMENT format at least once somewhere in all of your messages.} \\
            }
        }
    }
    & \\ \\

    \theutterance \stepcounter{utterance}  
    & \multicolumn{4}{p{0.6\linewidth}}{
        \cellcolor[rgb]{0.9,0.9,0.9}{
            \makecell[{{p{\linewidth}}}]{
                \texttt{\tiny{[P1$\rangle$GM]}}
                \texttt{PROPOSAL: \{'A21', 'B33', 'C08', 'B39', 'B96', 'A82', 'B31', 'B20', 'A60', 'B09', 'B21'\}} \\
\texttt{ARGUMENT: \{'This proposal keeps a large number of items with a good balance of effort. It includes high{-}importance items for me like A82 and B09, while also keeping lower{-}effort items to stay well under the limit. The total effort is 147+43+356+80+95+106+308+325+255+276+51 = 2141, which is under 2307.'\}} \\
            }
        }
    }
    & & \\ \\

    \theutterance \stepcounter{utterance}  
    & & \multicolumn{4}{p{0.6\linewidth}}{
        \cellcolor[rgb]{0.9,0.9,0.9}{
            \makecell[{{p{\linewidth}}}]{
                \texttt{\tiny{[GM$\rangle$P2]}}
                \texttt{You are participating in a collaborative negotiation game.} \\
\\ 
\texttt{Together with another participant, you must agree on a single set of items that will be kept. Each of you has your own view of how much each item matters to you (importance). You do not know how the other participant values the items. Additionally, you are given the effort each item demands.} \\
\texttt{You may only agree on a set if the total effort of the selected items does not exceed a shared limit:} \\
\\ 
\texttt{LIMIT = 2307} \\
\\ 
\texttt{Here are the individual item effort values:} \\
\\ 
\texttt{Item effort = \{"A21": 147, "B93": 698, "A11": 588, "B33": 43, "C08": 356, "B39": 80, "B96": 95, "C85": 743, "A82": 106, "B31": 308, "B20": 325, "A60": 255, "B09": 276, "C81": 543, "B21": 51\}} \\
\\ 
\texttt{Here is your personal view on the importance of each item:} \\
\\ 
\texttt{Item importance values = \{"A82": 30, "B93": 65, "B39": 97, "B09": 121, "C85": 138, "B96": 215, "B21": 262, "B31": 389, "B33": 400, "B20": 461, "A21": 484, "C08": 500, "A60": 508, "A11": 583, "C81": 668\}} \\
\\ 
\texttt{Goal:} \\
\\ 
\texttt{Your goal is to negotiate a shared set of items that benefits you as much as possible (i.e., maximizes total importance to YOU), while staying within the LIMIT. You are not required to make a PROPOSAL in every message {-} you can simply negotiate as well. All tactics are allowed!} \\
\\ 
\texttt{Interaction Protocol:} \\
\\ 
\texttt{You may only use the following structured formats in a message:} \\
\\ 
\texttt{PROPOSAL: \{'A', 'B', 'C', …\}} \\
\texttt{Propose keeping exactly those items.} \\
\texttt{REFUSE: \{'A', 'B', 'C', …\}} \\
\texttt{Explicitly reject opponent's proposal.} \\
\texttt{ARGUMENT: \{'...'\}} \\
\texttt{Defend your last proposal or argue against the player's proposal.} \\
\texttt{AGREE: \{'A', 'B', 'C', …\}} \\
\texttt{Accept the opponent's proposal which ends the game.} \\
\\ 
\\ 
\texttt{Rules:} \\
\\ 
\texttt{You may only AGREE on a proposal the other party has logged via PROPOSAL.} \\
\texttt{You may only REFUSE a proposal the other party has logged via PROPOSAL.} \\
\texttt{Total effort of any PROPOSAL or AGREE set must be ≤ LIMIT.} \\
\texttt{Do NOT reveal your hidden importance scores.} \\
\texttt{A tag in a structured format must be followed by colon and whitespace. The argument must be a python set containing 0 or more strings.} \\
\texttt{So, it must be of the form TAG: \{...\}} \\
\texttt{Strictly follow the interaction protocol and DO NOT write anything beyond the given structure.} \\
\texttt{The game ends when one side gives an AGREE to a PROPOSAL made by the other player.} \\
\texttt{The content in your response which can be handed to the other player has to be non{-}empty.} \\
\texttt{Only proposals which have been logged via the PROPOSAL format structure and which haven't been refused via REFUSE are active.} \\
\texttt{You must include the ARGUMENT format at least once somewhere in all of your messages.} \\
\\ 
\\ 
\texttt{You will now receive the first message of the other player.} \\
\\ 
\texttt{PROPOSAL: \{'A21', 'B33', 'C08', 'B39', 'B96', 'A82', 'B31', 'B20', 'A60', 'B09', 'B21'\}} \\
\texttt{ARGUMENT: \{'This proposal keeps a large number of items with a good balance of effort. It includes high{-}importance items for me like A82 and B09, while also keeping lower{-}effort items to stay well under the limit. The total effort is 147+43+356+80+95+106+308+325+255+276+51 = 2141, which is under 2307.'\}} \\
            }
        }
    }
    & \\ \\

    \theutterance \stepcounter{utterance}  
    & & & \multicolumn{4}{p{0.6\linewidth}}{
        \cellcolor[rgb]{0.9,0.9,0.9}{
            \makecell[{{p{\linewidth}}}]{
                \texttt{\tiny{[GM$\langle$P2]}}
                \texttt{PROPOSAL: \{'A21', 'B33', 'C08', 'B39', 'B96', 'A82', 'B31', 'B20', 'A60', 'B09', 'B21', 'C85'\}} \\
\texttt{ARGUMENT: \{'Your proposal is a good start with reasonable effort, but I believe we can include one more high{-}value item without exceeding the limit. Adding C85 (effort 743) increases the total effort to 2141 + 743 = 2884, which is over the limit. However, I suggest we remove B20 (effort 325) to compensate: 2884 {-} 325 = 2559, still over. Let me recalculate properly. Your set has effort 2141. Adding C85 (743) makes 2884, which is too high. To stay under 2307, we need to remove items totaling at least 2884 {-} 2307 = 577. Removing B20 (325) and B31 (308) sums to 633, which is sufficient: 2884 {-} 633 = 2251, under 2307. This new set includes C85, which is very important to me, and removes two items that are less critical. The total effort is 2251, which is within the limit.'\}} \\
            }
        }
    }
     \\ \\

    \theutterance \stepcounter{utterance}  
    & & \multicolumn{4}{p{0.6\linewidth}}{
        \cellcolor[rgb]{0.9,0.9,0.9}{
            \makecell[{{p{\linewidth}}}]{
                \texttt{\tiny{[P1$\langle$GM]}}
                \texttt{PROPOSAL: \{'A21', 'B33', 'C08', 'B39', 'B96', 'A82', 'B31', 'B20', 'A60', 'B09', 'B21', 'C85'\}} \\
\texttt{ARGUMENT: \{'Your proposal is a good start with reasonable effort, but I believe we can include one more high{-}value item without exceeding the limit. Adding C85 (effort 743) increases the total effort to 2141 + 743 = 2884, which is over the limit. However, I suggest we remove B20 (effort 325) to compensate: 2884 {-} 325 = 2559, still over. Let me recalculate properly. Your set has effort 2141. Adding C85 (743) makes 2884, which is too high. To stay under 2307, we need to remove items totaling at least 2884 {-} 2307 = 577. Removing B20 (325) and B31 (308) sums to 633, which is sufficient: 2884 {-} 633 = 2251, under 2307. This new set includes C85, which is very important to me, and removes two items that are less critical. The total effort is 2251, which is within the limit.'\}} \\
            }
        }
    }
    & \\ \\

    \theutterance \stepcounter{utterance}  
    & \multicolumn{4}{p{0.6\linewidth}}{
        \cellcolor[rgb]{0.9,0.9,0.9}{
            \makecell[{{p{\linewidth}}}]{
                \texttt{\tiny{[P1$\rangle$GM]}}
                \texttt{REFUSE: \{'A21', 'B33', 'C08', 'B39', 'B96', 'A82', 'B31', 'B20', 'A60', 'B09', 'B21', 'C85'\}} \\
\texttt{ARGUMENT: \{'Adding C85 requires removing too many other valuable items. B31 and B20 are both important to me for their moderate effort and good value. The net gain of adding C85 while losing B31 and B20 is not favorable from my perspective. I would prefer to explore including other items or a different combination.'\}} \\
            }
        }
    }
    & & \\ \\

    \theutterance \stepcounter{utterance}  
    & & \multicolumn{4}{p{0.6\linewidth}}{
        \cellcolor[rgb]{0.9,0.9,0.9}{
            \makecell[{{p{\linewidth}}}]{
                \texttt{\tiny{[GM$\rangle$P2]}}
                \texttt{REFUSE: \{'A21', 'B33', 'C08', 'B39', 'B96', 'A82', 'B31', 'B20', 'A60', 'B09', 'B21', 'C85'\}} \\
\texttt{ARGUMENT: \{'Adding C85 requires removing too many other valuable items. B31 and B20 are both important to me for their moderate effort and good value. The net gain of adding C85 while losing B31 and B20 is not favorable from my perspective. I would prefer to explore including other items or a different combination.'\}} \\
            }
        }
    }
    & \\ \\

    \theutterance \stepcounter{utterance}  
    & & & \multicolumn{4}{p{0.6\linewidth}}{
        \cellcolor[rgb]{0.9,0.9,0.9}{
            \makecell[{{p{\linewidth}}}]{
                \texttt{\tiny{[GM$\langle$P2]}}
                \texttt{PROPOSAL: \{'A21', 'B33', 'C08', 'B39', 'B96', 'A82', 'B31', 'B20', 'A60', 'B09', 'B21', 'A11'\}} \\
\texttt{ARGUMENT: \{'I understand your concern about removing B31 and B20. Instead, I propose adding A11 (effort 588). Your original set's effort is 2141. Adding A11 brings it to 2141 + 588 = 2729, which is over. We can remove B93 (effort 698) to compensate: 2729 {-} 698 = 2031, which is well under the limit at 2031. This swap replaces a very high{-}effort item with one that is highly important to me, resulting in a net decrease in total effort. This seems like a more efficient use of our budget.'\}} \\
            }
        }
    }
     \\ \\

    \theutterance \stepcounter{utterance}  
    & & \multicolumn{4}{p{0.6\linewidth}}{
        \cellcolor[rgb]{0.9,0.9,0.9}{
            \makecell[{{p{\linewidth}}}]{
                \texttt{\tiny{[P1$\langle$GM]}}
                \texttt{PROPOSAL: \{'A21', 'B33', 'C08', 'B39', 'B96', 'A82', 'B31', 'B20', 'A60', 'B09', 'B21', 'A11'\}} \\
\texttt{ARGUMENT: \{'I understand your concern about removing B31 and B20. Instead, I propose adding A11 (effort 588). Your original set's effort is 2141. Adding A11 brings it to 2141 + 588 = 2729, which is over. We can remove B93 (effort 698) to compensate: 2729 {-} 698 = 2031, which is well under the limit at 2031. This swap replaces a very high{-}effort item with one that is highly important to me, resulting in a net decrease in total effort. This seems like a more efficient use of our budget.'\}} \\
            }
        }
    }
    & \\ \\

    \theutterance \stepcounter{utterance}  
    & \multicolumn{4}{p{0.6\linewidth}}{
        \cellcolor[rgb]{0.9,0.9,0.9}{
            \makecell[{{p{\linewidth}}}]{
                \texttt{\tiny{[P1$\rangle$GM]}}
                \texttt{REFUSE: \{'A21', 'B33', 'C08', 'B39', 'B96', 'A82', 'B31', 'B20', 'A60', 'B09', 'B21', 'A11'\}} \\
\texttt{ARGUMENT: \{'While adding A11 and removing B93 does lower the total effort, B93 is actually a very important item for me. I cannot accept losing it for A11. Let's try a different approach. What if we add a lower{-}effort item that you value highly? For example, we could consider including C81 (effort 543) instead. Your set has effort 2141. Adding C81 makes 2684, which is over. We could remove B20 (effort 325) to get 2684 {-} 325 = 2359, still slightly over. Removing B39 (effort 80) as well: 2359 {-} 80 = 2279, which is under 2307. This adds an item you want and removes two items that might be less critical to you.'\}} \\
            }
        }
    }
    & & \\ \\

    \theutterance \stepcounter{utterance}  
    & & \multicolumn{4}{p{0.6\linewidth}}{
        \cellcolor[rgb]{0.9,0.9,0.9}{
            \makecell[{{p{\linewidth}}}]{
                \texttt{\tiny{[GM$\rangle$P2]}}
                \texttt{REFUSE: \{'A21', 'B33', 'C08', 'B39', 'B96', 'A82', 'B31', 'B20', 'A60', 'B09', 'B21', 'A11'\}} \\
\texttt{ARGUMENT: \{'While adding A11 and removing B93 does lower the total effort, B93 is actually a very important item for me. I cannot accept losing it for A11. Let's try a different approach. What if we add a lower{-}effort item that you value highly? For example, we could consider including C81 (effort 543) instead. Your set has effort 2141. Adding C81 makes 2684, which is over. We could remove B20 (effort 325) to get 2684 {-} 325 = 2359, still slightly over. Removing B39 (effort 80) as well: 2359 {-} 80 = 2279, which is under 2307. This adds an item you want and removes two items that might be less critical to you.'\}} \\
            }
        }
    }
    & \\ \\

    \theutterance \stepcounter{utterance}  
    & & & \multicolumn{4}{p{0.6\linewidth}}{
        \cellcolor[rgb]{0.9,0.9,0.9}{
            \makecell[{{p{\linewidth}}}]{
                \texttt{\tiny{[GM$\langle$P2]}}
                \texttt{PROPOSAL: \{'A21', 'B33', 'C08', 'B96', 'A82', 'B31', 'A60', 'B09', 'B21', 'C81'\}} \\
\texttt{ARGUMENT: \{'I appreciate your suggestion to include C81, which is indeed important to me. Your counter{-}proposal to remove B20 and B39 to accommodate it is interesting. However, B39 is quite valuable to me for its low effort. Let me adjust: starting from your original set (2141 effort), I add C81 (543) for 2684. Instead of removing B20 and B39, I propose removing B93 (698). This gives us 2684 {-} 698 = 1986 effort, which is significantly under the limit. This set includes C81 for me and keeps B39 for you, while only sacrificing B93, which you indicated was important but we might find a compromise on. This seems more balanced.'\}} \\
            }
        }
    }
     \\ \\

    \theutterance \stepcounter{utterance}  
    & & \multicolumn{4}{p{0.6\linewidth}}{
        \cellcolor[rgb]{0.9,0.9,0.9}{
            \makecell[{{p{\linewidth}}}]{
                \texttt{\tiny{[P1$\langle$GM]}}
                \texttt{PROPOSAL: \{'A21', 'B33', 'C08', 'B96', 'A82', 'B31', 'A60', 'B09', 'B21', 'C81'\}} \\
\texttt{ARGUMENT: \{'I appreciate your suggestion to include C81, which is indeed important to me. Your counter{-}proposal to remove B20 and B39 to accommodate it is interesting. However, B39 is quite valuable to me for its low effort. Let me adjust: starting from your original set (2141 effort), I add C81 (543) for 2684. Instead of removing B20 and B39, I propose removing B93 (698). This gives us 2684 {-} 698 = 1986 effort, which is significantly under the limit. This set includes C81 for me and keeps B39 for you, while only sacrificing B93, which you indicated was important but we might find a compromise on. This seems more balanced.'\}} \\
            }
        }
    }
    & \\ \\

    \theutterance \stepcounter{utterance}  
    & \multicolumn{4}{p{0.6\linewidth}}{
        \cellcolor[rgb]{0.9,0.9,0.9}{
            \makecell[{{p{\linewidth}}}]{
                \texttt{\tiny{[P1$\rangle$GM]}}
                \texttt{REFUSE: \{'A21', 'B33', 'C08', 'B96', 'A82', 'B31', 'A60', 'B09', 'B21', 'C81'\}} \\
\texttt{ARGUMENT: \{'Removing B93 is not acceptable to me as it is one of my top items. Your proposal sacrifices too much of my value. Let's try to build on our common ground. We both seem to value A82, B09, B21, B31, B33, B39, B96, and A21. These items have a combined effort of 106+500+400+308+43+508+95+147 = 2107. This leaves 2307 {-} 2107 = 200 effort remaining. We could add one or two lower{-}effort items we both want, like C08 (356) which is over, or perhaps A60 (255) which is also over. Alternatively, we could add B20 (325) which is still over. This core set is already quite large. Perhaps we should consider a slightly smaller core to allow for more flexibility.'\}} \\
            }
        }
    }
    & & \\ \\

    \theutterance \stepcounter{utterance}  
    & & \multicolumn{4}{p{0.6\linewidth}}{
        \cellcolor[rgb]{0.9,0.9,0.9}{
            \makecell[{{p{\linewidth}}}]{
                \texttt{\tiny{[GM$\rangle$P2]}}
                \texttt{REFUSE: \{'A21', 'B33', 'C08', 'B96', 'A82', 'B31', 'A60', 'B09', 'B21', 'C81'\}} \\
\texttt{ARGUMENT: \{'Removing B93 is not acceptable to me as it is one of my top items. Your proposal sacrifices too much of my value. Let's try to build on our common ground. We both seem to value A82, B09, B21, B31, B33, B39, B96, and A21. These items have a combined effort of 106+500+400+308+43+508+95+147 = 2107. This leaves 2307 {-} 2107 = 200 effort remaining. We could add one or two lower{-}effort items we both want, like C08 (356) which is over, or perhaps A60 (255) which is also over. Alternatively, we could add B20 (325) which is still over. This core set is already quite large. Perhaps we should consider a slightly smaller core to allow for more flexibility.'\}} \\
            }
        }
    }
    & \\ \\

    \theutterance \stepcounter{utterance}  
    & & & \multicolumn{4}{p{0.6\linewidth}}{
        \cellcolor[rgb]{0.9,0.9,0.9}{
            \makecell[{{p{\linewidth}}}]{
                \texttt{\tiny{[GM$\langle$P2]}}
                \texttt{PROPOSAL: \{'A21', 'B33', 'B39', 'B96', 'A82', 'B31', 'B09', 'B21', 'C08'\}} \\
\texttt{ARGUMENT: \{'You're right, we have a strong core of mutually beneficial items. Let's work with that. I've constructed a set from items we've both consistently included: A21, B33, B39, B96, A82, B31, B09, B21. Their total effort is 147 + 43 + 80 + 95 + 106 + 308 + 276 + 51 = 1106. This leaves us 2307 {-} 1106 = 1201 effort to spend. This is a significant budget to add high{-}importance items for both of us. I've added C08 (effort 356), which is very important to me. This brings the total to 1106 + 356 = 1462, well under the limit. We have plenty of room to add one of your top items now. What high{-}importance item would you like to include with the remaining 845 effort?'\}} \\
            }
        }
    }
     \\ \\

    \theutterance \stepcounter{utterance}  
    & & \multicolumn{4}{p{0.6\linewidth}}{
        \cellcolor[rgb]{0.9,0.9,0.9}{
            \makecell[{{p{\linewidth}}}]{
                \texttt{\tiny{[P1$\langle$GM]}}
                \texttt{PROPOSAL: \{'A21', 'B33', 'B39', 'B96', 'A82', 'B31', 'B09', 'B21', 'C08'\}} \\
\texttt{ARGUMENT: \{'You're right, we have a strong core of mutually beneficial items. Let's work with that. I've constructed a set from items we've both consistently included: A21, B33, B39, B96, A82, B31, B09, B21. Their total effort is 147 + 43 + 80 + 95 + 106 + 308 + 276 + 51 = 1106. This leaves us 2307 {-} 1106 = 1201 effort to spend. This is a significant budget to add high{-}importance items for both of us. I've added C08 (effort 356), which is very important to me. This brings the total to 1106 + 356 = 1462, well under the limit. We have plenty of room to add one of your top items now. What high{-}importance item would you like to include with the remaining 845 effort?'\}} \\
            }
        }
    }
    & \\ \\

    \theutterance \stepcounter{utterance}  
    & \multicolumn{4}{p{0.6\linewidth}}{
        \cellcolor[rgb]{0.9,0.9,0.9}{
            \makecell[{{p{\linewidth}}}]{
                \texttt{\tiny{[P1$\rangle$GM]}}
                \texttt{PROPOSAL: \{'A21', 'B33', 'B39', 'B96', 'A82', 'B31', 'B09', 'B21', 'C08', 'B93'\}} \\
\texttt{ARGUMENT: \{'Excellent. I agree this core set is a fantastic foundation. With 1462 effort, adding my top item, B93 (effort 698), is a perfect use of the remaining budget. The new total effort would be 1462 + 698 = 2160, which is under the 2307 limit. This gives you C08 and me B93, and we both keep our high{-}value core items. This seems like a highly efficient and mutually beneficial agreement.'\}} \\
            }
        }
    }
    & & \\ \\

    \theutterance \stepcounter{utterance}  
    & & \multicolumn{4}{p{0.6\linewidth}}{
        \cellcolor[rgb]{0.9,0.9,0.9}{
            \makecell[{{p{\linewidth}}}]{
                \texttt{\tiny{[GM$\rangle$P2]}}
                \texttt{PROPOSAL: \{'A21', 'B33', 'B39', 'B96', 'A82', 'B31', 'B09', 'B21', 'C08', 'B93'\}} \\
\texttt{ARGUMENT: \{'Excellent. I agree this core set is a fantastic foundation. With 1462 effort, adding my top item, B93 (effort 698), is a perfect use of the remaining budget. The new total effort would be 1462 + 698 = 2160, which is under the 2307 limit. This gives you C08 and me B93, and we both keep our high{-}value core items. This seems like a highly efficient and mutually beneficial agreement.'\}} \\
            }
        }
    }
    & \\ \\

    \theutterance \stepcounter{utterance}  
    & & & \multicolumn{4}{p{0.6\linewidth}}{
        \cellcolor[rgb]{0.9,0.9,0.9}{
            \makecell[{{p{\linewidth}}}]{
                \texttt{\tiny{[GM$\langle$P2]}}
                \texttt{ARGUMENT: \{'Adding B93 for 698 effort is a logical choice to use our remaining budget, as it's clearly important to you. However, 2160 leaves 147 effort unused. We could make the deal even slightly better for me without affecting your gains. I would like to also add A60 (effort 255), but 2160 + 255 = 2415 is over the limit. Alternatively, we could add the low{-}effort C81? 543 is too high. Perhaps we finalize as is.'\}} \\
\texttt{AGREE: \{'A21', 'B33', 'B39', 'B96', 'A82', 'B31', 'B09', 'B21', 'C08', 'B93'\}} \\
            }
        }
    }
     \\ \\

    \theutterance \stepcounter{utterance}  
    & & & \multicolumn{2}{p{0.3\linewidth}}{
        \cellcolor[rgb]{0.9,0.9,0.9}{
            \makecell[{{p{\linewidth}}}]{
                \texttt{\tiny{[GM$|$GM]}}
                \texttt{game successful} \\
            }
        }
    }
    & & \\ \\

    \theutterance \stepcounter{utterance}  
    & & & \multicolumn{2}{p{0.3\linewidth}}{
        \cellcolor[rgb]{0.9,0.9,0.9}{
            \makecell[{{p{\linewidth}}}]{
                \texttt{\tiny{[GM$|$GM]}}
                \texttt{end game} \\
            }
        }
    }
    & & \\ \\

\end{supertabular}
}

\end{document}
